\section{Mato Grosso: Contexto e Dados}
\label{sec:dados}

\begingroup
\clubpenalty=10000
\widowpenalty=10000
\input{../stats/descritiva/resultados_h1/tabelas/numeros}

O recorte descritivo reúne informações de 2001 a 2024 para uma grade de 141 municípios de Mato Grosso, definida pela malha do IBGE de 2023. Transições de cobertura vegetal, produção agropecuária, preços, capacidade produtiva e registros de autuação permitem examinar dimensões distintas do problema. A cobertura efetivamente observada varia entre fontes, por isso, a grade de \Hobs{} pares município-ano não constitui uma amostra completa para todas as medidas.

\subsection{Expansão agropecuária e transformação da cobertura vegetal}
\label{subsec:mt-expansao}

A expansão agropecuária e a transformação da cobertura vegetal constituem duas dimensões relacionadas, mas não equivalentes, da alocação econômica da terra. A discussão sobre retornos privados e custos públicos, desenvolvida na Seção~\ref{sec:literatura}, situa a conversão na comparação entre rendas agropecuárias e benefícios da vegetação que nem sempre são apropriados por quem decide seu uso \cite{Angelsen2010}. Para Mato Grosso, cabe estabelecer quanto cresceu a atividade produtiva e como evoluiu a conversão observada, sem deduzir da primeira a magnitude da segunda.

A Figura~\ref{fig:h1-lavouras} apresenta área plantada, área colhida e produção de soja e milho segundo a Produção Agrícola Municipal (PAM/IBGE). A amostra disponível reúne os municípios com as três medidas observadas em cada ano; a fixa exige essa informação nos 24 anos, separadamente por cultura. Nos 66 municípios fixos da soja, a produção passou de 9,51 para 30,67 milhões de toneladas entre 2001 e 2024, enquanto a área colhida aumentou de 3,11 para 9,88 milhões de hectares. Nos 121 municípios fixos do milho, a produção passou de 1,72 para 45,03 milhões de toneladas, e a área colhida, de 0,53 para 6,70 milhões de hectares. São totais amostrais, distintos dos totais estaduais e do intervalo 2000--2024 apresentados na revisão.

\begin{figure}[htbp]
    \centering
    \includegraphics[width=\textwidth]{../stats/figuras/fig05_lavouras.pdf}
    \caption{Expansão das lavouras de soja e milho nas amostras disponível e fixa}
    \label{fig:h1-lavouras}
    \begin{minipage}{\textwidth}\footnotesize
    Fonte: elaboração própria com IBGE/PAM. Área plantada e colhida em milhões de hectares; produção em milhões de toneladas. A amostra fixa contém 66 municípios para soja e 121 para milho, com plantio, colheita e produção observados em todos os anos. Ausências não são zeros. As somas de áreas de culturas diferentes não representam necessariamente hectares fisicamente distintos. A geografia municipal histórica das séries administrativas não foi reconstituída.
    \end{minipage}
\end{figure}

O crescimento da produção combina mudanças de área e de rendimento, com diferenças entre atividades. Calculado pela razão entre produção total e área colhida total dos mesmos municípios, o rendimento do milho passou de 3,26 para 6,72 toneladas por hectare; o da soja, de 3,06 para 3,10. Essa comparação de pontas não resume as oscilações intermediárias nem permite concluir que a soja não tenha incorporado mudanças tecnológicas. A possibilidade de duas safras na mesma área, examinada por \citeonline{VanWeyEtAl2013}, ajuda a compreender por que o avanço das áreas declaradas de soja e milho não pode ser somado como expansão física da fronteira. A PAM não identifica as parcelas com cultivos sucessivos ou abertas sobre vegetação nativa.

A escala da pecuária e da comercialização externa complementa esse quadro (Tabela~\ref{tab:h1-oferta-contexto}). O rebanho dos 139 municípios com observação em todo o período cresceu 64,5\%, enquanto o número estadual de bovinos abatidos sob inspeção aumentou 246,9\%. A diferença reforça a necessidade de distinguir estoque de animais e fluxo de abate: nenhum deles mede diretamente a extensão das pastagens. O peso das carcaças descreve outra dimensão da oferta, sem equivaler ao volume de carne desossada. Os valores exportados indicam ampliação das operações registradas por domicílio fiscal, mas seu crescimento nominal combina mudanças de quantidade e preço e não mede o lucro dos produtores. A base inicial muito pequena das exportações de milho recomenda cautela adicional na leitura de sua expansão.

\input{../stats/tabelas/tab_contexto_oferta_h1}
\FloatBarrier

Para observar a transformação da cobertura vegetal, utilizam-se as transições anuais do MapBiomas, Coleção 10.1: vegetação nativa em $t-1$ convertida para agricultura ou pastagem em $t$. A taxa divide a área convertida pelo estoque nativo inicial na mesma máscara espacial. Trata-se de conversão bruta por destino, distinta da perda líquida de cobertura ou da diferença entre estoques de área agrícola. O indicador não registra como nova conversão a passagem de pastagem já existente para lavoura, nem identifica isoladamente a causa ou a regularidade jurídica da mudança entre mapas.

Essa medida está disponível em \Hvalid{} pares município-ano. Os \Hmissing{} restantes, ou \HmissingPct{}\% da grade, permanecem ausentes pela exigência de cobertura integral do par de mapas; essa proporção de pares não equivale à proporção de área territorial sem observação. A Figura~\ref{fig:h1-conversao} compara os municípios disponíveis em cada ano com os \Hfixed{} observados em todo o período. Na amostra fixa, a taxa agregada para agricultura passou de \HagFirst{}\% em 2001 para \HagLast{}\% em 2024; para pastagem, de \HpaFirst{}\% para \HpaLast{}\%. A trajetória não foi monotônica, sobretudo para pastagem. Considerando a soma das conversões observadas nos dois destinos em todos os pares válidos, a pastagem respondeu por \HpastShare{}\%; essa participação não é uma estimativa da composição da conversão total do estado.

\begin{figure}[htbp]
    \centering
    \includegraphics[width=\textwidth]{../stats/descritiva/resultados_h1/figuras/fig01_conversao.pdf}
    \caption{Taxas agregadas de conversão da vegetação nativa para agricultura e pastagem}
    \label{fig:h1-conversao}
    \begin{minipage}{\textwidth}\footnotesize
    Fonte: elaboração própria com MapBiomas 10.1 e malha IBGE 2023. Cada ponto corresponde à soma da conversão dividida pela soma da vegetação nativa inicial dos mesmos municípios, multiplicada por 100. A amostra disponível varia entre \HnMin{} e \HnMax{} municípios por ano; a fixa contém \Hfixed{}. Pares incompletos não entram no numerador nem no denominador. A razão entre somas atribui maior peso aos municípios com maior estoque nativo inicial; não é a média simples das taxas municipais.
    \end{minipage}
\end{figure}

\FloatBarrier
O crescimento produtivo e as taxas menores de conversão nos extremos do período qualificam uma leitura que vincule automaticamente aumento da oferta à abertura de novas áreas. O contraste dialoga com \citeonline{MacedoEtAl2012}, que documentaram crescimento da soja e redução do desmatamento na porção florestal de Mato Grosso, com expansão em terras anteriormente abertas. Os resultados aqui descritos abrangem outro período e amostras distintas; não identificam quanto da produção adicional decorreu da substituição de pastagens, da intensificação ou de novas conversões. Tampouco quantificam lucros privados, serviços ecossistêmicos perdidos ou perdas de bem-estar. Acrescentam à discussão sobre custos públicos a identificação dos fluxos que transformam a vegetação, mantendo em aberto as condições econômicas e institucionais que os orientaram.

\FloatBarrier
\subsection{Preços agropecuários e incentivos à conversão}
\label{subsec:mt-precos}

Os preços agropecuários representam uma dimensão observável dos incentivos à conversão discutidos na subseção correspondente da revisão. Mantidas as demais condições, preços mais altos podem elevar a receita esperada por hectare, com intensidade dependente da capacidade produtiva local. A evidência de \citeonline{BermanEtAl2023}, que combinaram preços internacionais e aptidão agrícola para examinar o desmatamento tropical, torna pertinente investigar essa exposição diferenciada. A previsão, entretanto, depende de custos, alternativas de uso e instituições; a cotação do produto não constitui medida completa de rentabilidade.

As séries da Figura~\ref{fig:h1-precos} provêm da \textit{Pink Sheet} do Banco Mundial e foram deflacionadas mês a mês pelo CPI-U dos Estados Unidos, disponibilizado pelo BLS/FRED, para dólares de poder de compra médio de 2024. Soja e milho são cotados por tonelada, e carne bovina, por quilograma. Cada média anual exige doze meses; o preço do ano civil anterior é associado à conversão em $t$, de modo que a janela representada é 2000--2023. O índice com base 100 em 2000 permite comparar movimentos proporcionais entre produtos, sem equiparar suas unidades ou receitas por hectare.

\begin{figure}[htbp]
    \centering
    \includegraphics[width=\textwidth]{../stats/descritiva/resultados_h1/figuras/fig03_precos.pdf}
    \caption{Trajetórias dos preços reais internacionais de soja, milho e carne bovina}
    \label{fig:h1-precos}
    \begin{minipage}{\textwidth}\footnotesize
    Fonte: elaboração própria com Banco Mundial e CPI-U BLS/FRED. Índice do preço real anual de cada produto, com 2000 = 100. São 24 valores anuais distintos por produto, repetidos na grade municipal; a repetição não amplia a informação temporal. Preços de $t-1$ são associados às transições de $t$. Dólares constantes não equivalem a preços locais em reais.
    \end{minipage}
\end{figure}

Em 2012, os índices alcançaram 210,7 para soja e 252,6 para milho, ante 163,8 para carne bovina. A valorização dos grãos não foi acompanhada de maneira proporcional pela carne, cuja trajetória diferiu nos anos seguintes. Essa distinção aproxima a leitura dos dados da questão examinada por \citeonline{Braganca2018}: no Tapajós, aumentos do preço relativo de lavouras em relação à carne estiveram associados à substituição de pastagens por cultivos e à redução do desmatamento local. Aqui, as séries de preços evidenciam movimentos relativos, mas as medidas de conversão selecionadas não observam a substituição entre usos já agropecuários. Não permitem, portanto, atribuir a trajetória de Mato Grosso ao mecanismo estimado naquele estudo.

Há outro contraste relevante: em 2023, os preços reais dos três produtos permaneciam acima dos valores de 2000, enquanto as taxas de conversão de 2024 eram inferiores às de 2001 na amostra fixa. Essa coexistência afasta uma leitura mecânica segundo a qual preços mais altos deveriam acompanhar taxas agregadas sempre maiores. Não constitui rejeição de H1, cuja formulação é condicional à capacidade produtiva e ao ambiente institucional, nem teste comparável aos desenhos empíricos da literatura. As séries internacionais não incorporam diretamente câmbio, custos locais, crédito ou expectativas individuais; sua defasagem organiza a informação, mas não assegura exogeneidade. A questão passa a ser em quais condições territoriais e institucionais esses movimentos se associam à conversão, e não se duas tendências agregadas coincidem.

\FloatBarrier
\subsection{Instituições ambientais e respostas aos incentivos}
\label{subsec:mt-instituicoes}

As instituições ambientais podem modificar os custos e os riscos associados às respostas aos incentivos de mercado. Na discussão sobre instituições e respostas aos incentivos da Seção~\ref{sec:literatura}, a distinção entre regras e condições de implementação impede tomar a existência de uma norma como medida de seu cumprimento. A evidência de \citeonline{AssuncaoEtAl2023}, baseada nas limitações de detecção por nuvens do DETER, identificou contenção do desmatamento pelo monitoramento e pela fiscalização na Amazônia. Para dialogar com esse resultado, os registros de Mato Grosso precisam ser examinados como evidência administrativa, sem lhes atribuir automaticamente a capacidade de medir efetividade.

O indicador disponível é o número de autos de infração do IBAMA classificados como Flora, por ano de lavratura e município atribuído pelas coordenadas à malha de referência. A categoria abrange infrações além da conversão, como transporte de produtos florestais. A contagem não representa o valor de multas aplicadas, o montante pago ou a confirmação definitiva de cada infração. A Figura~\ref{fig:h1-ibama} apresenta conjuntamente sua cobertura e as estatísticas das contagens com classificação temática completa nos registros utilizados.

\begin{figure}[htbp]
    \centering
    \includegraphics[width=\textwidth]{../stats/descritiva/resultados_h1/figuras/fig04_ibama.pdf}
    \caption{Cobertura e distribuição das autuações do IBAMA classificadas como Flora}
    \label{fig:h1-ibama}
    \begin{minipage}{\textwidth}\footnotesize
    Fonte: elaboração própria com registros do IBAMA e malha IBGE 2023. Painel A: municípios com contagem observada e com classificação temática completa nos registros. Painel B: média e mediana entre os municípios do segundo conjunto, de composição variável. O período contém \HibValid{} contagens observadas e \HibComplete{} com classificação completa. Ausência não é zero; classificação completa não certifica a completude da fonte ou da fiscalização. Os painéis possuem unidades e escalas distintas.
    \end{minipage}
\end{figure}

O número de municípios com classificação completa variou entre \HibMin{} e \HibMax{} por ano, o que limita a comparação direta dos níveis das contagens. Em 2018, a média foi de 18,77 autos e a mediana, de 5,5; em 2024, foram 4,41 e 3, respectivamente. Nos dois anos havia 82 municípios com classificação completa, mas a igualdade de tamanho não assegura a mesma composição. A distância entre média e mediana indica que as contagens mais elevadas em parte dos municípios influenciam o resumo agregado. Os zeros identificam ausência de autos de Flora entre os registros integralmente classificados daquele município-ano; não comprovam ausência de infrações ou de atuação fiscalizatória.

A redução da média entre esses anos não permite concluir que a fiscalização se tornou mais eficaz ou menos presente. As autuações podem aumentar por maior incidência de infrações, maior esforço de fiscalização ou melhor registro e localização, e esses processos podem ocorrer simultaneamente. Pela mesma razão, eventual associação positiva com conversão não demonstraria ineficácia institucional. O contraste com o estudo de \citeonline{AssuncaoEtAl2023} reside justamente no alcance da evidência: suas condições de identificação não são reproduzidas pela descrição de contagens. Os dados observam uma dimensão da atuação administrativa federal, deixando fora medidas abrangentes de cumprimento, normas informais, ação de outras organizações e governança territorial. Para H1, permanece necessário justificar em que sentido o indicador escolhido pode representar o condicionamento institucional da resposta aos incentivos.

\FloatBarrier
\subsection{Heterogeneidade econômica e ambiental de Mato Grosso}
\label{subsec:mt-heterogeneidade}

A heterogeneidade econômica e ambiental de Mato Grosso qualifica as tendências agregadas discutidas até aqui. Na subseção \emph{Situação econômica e ambiental de Mato Grosso}, a revisão destacou a diversidade de trajetórias documentada pelos mapas de \citeonline{SimoesEtAl2020}. O recorte municipal desta análise permite examinar como se distribuem conversão e capacidade produtiva, mas não reproduz a observação espacialmente explícita de transições dentro de cada bioma. Essa diferença de escala precisa anteceder a comparação territorial.

Cada município foi agrupado pelo bioma com a maior fração de sua área, obtida pela interseção da malha municipal de 2023 com os biomas do IBGE de 2019. As frações originais foram preservadas, e \Hmixed{} municípios apresentam mais de um bioma. A Tabela~\ref{tab:h1-cobertura} mostra que a cobertura da conversão também é desigual: 71,6\% dos pares potenciais no grupo amazônico, 96,4\% no Cerrado e 56,2\% no Pantanal. Dos quatro municípios de predominância pantaneira, somente um possui conversão observada em todos os anos. A amostra fixa reduz mudanças de composição temporal, mas essa seleção não produz grupos igualmente informativos nem resolve a compatibilidade histórica de todas as séries administrativas.

\input{../stats/descritiva/resultados_h1/tabelas/tab01_cobertura}

A Figura~\ref{fig:h1-distribuicao} dá peso igual aos anos ao calcular a média temporal das 24 taxas de cada município fixo, distinguindo esse objeto da razão entre somas apresentada na Figura~\ref{fig:h1-conversao}. Para pastagem, a mediana dessas médias foi de 1,42\% nos 39 municípios de predominância amazônica e de 0,73\% nos 55 de predominância Cerrado. Dentro dos dois grupos há dispersão expressiva e valores superiores aos resumos centrais. Para agricultura, as taxas médias são menores, mas igualmente heterogêneas. O único ponto pantaneiro não oferece base para caracterizar uma distribuição ou hierarquizar o bioma em relação aos demais.

\begin{figure}[htbp]
    \centering
    \includegraphics[width=\textwidth]{../stats/descritiva/resultados_h1/figuras/fig02_distribuicao.pdf}
    \caption{Heterogeneidade das taxas médias municipais de conversão, 2001--2024}
    \label{fig:h1-distribuicao}
    \begin{minipage}{\textwidth}\footnotesize
    Fonte: elaboração própria com MapBiomas 10.1 e IBGE. Cada ponto é a média aritmética das 24 taxas anuais de um município da amostra fixa. A posição horizontal dentro do grupo apenas separa os pontos; não representa localização geográfica. Grupos definidos pelo bioma predominante na área municipal, sem rateio da conversão pelas frações de bioma. Todos os municípios fixos são mostrados; o grupo pantaneiro contém uma única observação municipal.
    \end{minipage}
\end{figure}

Essas diferenças não isolam um efeito do bioma. Os municípios agrupados compartilham apenas um critério de predominância e podem diferir em vegetação disponível, estrutura produtiva, acesso a mercados e condições institucionais. Os fluxos não foram repartidos proporcionalmente pelas frações de bioma, pois isso atribuiria às transições uma localização que o agregado não informa. A evidência complementa a caracterização territorial da literatura ao mostrar que a variação dentro dos grupos também precisa ser considerada na investigação de respostas aos incentivos.

Outra dimensão dessa heterogeneidade aparece no rendimento potencial do GAEZ v4, apresentado na Tabela~\ref{tab:h1-capacidade}. No cenário de alto insumo e sequeiro, a mediana municipal da soja foi de 2.701 kg de matéria seca por hectare no grupo amazônico e de 3.261 no Cerrado; para milho, de 3.002 e 3.320, respectivamente. Os intervalos interquartis mostram que as possibilidades produtivas não se resumem à posição dessas medianas. A tabela abrange os 141 municípios, cada um considerado uma única vez, e não apenas os 95 da comparação de conversão. As posições relativas dos grupos nos dois resumos não constituem, assim, uma estimativa da relação entre capacidade e conversão.

\input{../stats/descritiva/resultados_h1/tabelas/tab02_capacidade}
\FloatBarrier

O rendimento potencial difere do rendimento efetivamente obtido nas áreas colhidas da PAM: este responde ao manejo, aos insumos e à seleção das parcelas cultivadas. No GAEZ, os pixels válidos foram agregados por município com ponderação pelo cosseno da latitude, aproximando diferenças de área; os resumos entre municípios lhes atribuem peso igual. A referência combina clima de 1961--1990 e informações de terras de 2010, não garantindo predeterminação integral em relação ao painel. Além disso, matéria seca e unidades comerciais precisam ser compatibilizadas antes de qualquer combinação com preços, e não há medida validada de capacidade pecuária neste conjunto. A heterogeneidade descrita torna pertinente o exame proposto em H1, mas a escolha de medidas predeterminadas e o suporte das comparações permanecem questões da etapa econométrica.

\subsection{Síntese e articulação com a estratégia econométrica}
\label{subsec:mt-sintese}

A descrição de Mato Grosso reúne crescimento agropecuário, taxas de conversão menores entre os extremos do período, predominância da pastagem nos fluxos observados, movimentos distintos de preços e registros institucionais de cobertura desigual. A dispersão municipal e as diferenças de capacidade produtiva mostram por que médias estaduais não esgotam o problema de como os retornos econômicos esperados e o ambiente institucional condicionam a taxa de conversão, nem como essa relação varia entre atividades e biomas. A construção de medidas brutas por destino atende ao primeiro objetivo específico e delimita o suporte dos objetivos de examinar preços em interação com capacidade predeterminada, condicionamento institucional e heterogeneidade. Mantém-se a formulação de H1: \emph{Elevações nos retornos econômicos esperados dos usos agropecuários aumentam a taxa de conversão da vegetação nativa, sendo essa resposta, \textit{ceteris paribus}, amplificada pela capacidade produtiva local da terra e modulada pela efetividade do ambiente institucional.} Os resultados reconhecem a coexistência dessas dimensões, sem confirmar ou rejeitar a hipótese; permanece em aberto se e em que intensidade os incentivos se associam à conversão sob diferentes condições locais.

Responder a essa questão requer distinguir a variação dos incentivos de tendências comuns, diferenças persistentes entre territórios, seleção amostral e respostas da fiscalização às próprias infrações. A passagem à Seção~\ref{sec:estrategia} envolve desenvolver o desenho geral anunciado na introdução: dados em painel, ordenação temporal entre preços e conversão, interações com capacidade produtiva e condições institucionais, efeitos fixos territoriais e temporais e controles compatíveis com a investigação. A especificação deverá justificar a mensuração, a amostra e o suporte de cada comparação, especialmente diante da cobertura restrita do Pantanal e das limitações temporais do GAEZ. Esses procedimentos não asseguram identificação causal; quando suas condições não puderem ser sustentadas, os resultados deverão ser interpretados como associações condicionais. A descrição estabelece, portanto, os objetos e os limites a partir dos quais a investigação econométrica de H1 poderá avançar.

\endgroup
